\section{Why This Document Exists --- and Five Ideas You Need First}\label{sec:intro}

In 2026, a quiet rule change in Brussels began deciding which small factories in
southern Taiwan keep their European customers. This proposal explains that
situation from zero, shows with verified numbers why it matters, and presents a
solution designed for Taiwan's 2026 Presidential Hackathon. Before the story,
five concepts.

\begin{concept}{Carbon pricing.}
Governments increasingly make emitters pay per tonne of carbon dioxide
equivalent (tCO$_2$e) released. Taiwan does this with a \emph{carbon fee} on
large emitters; the EU does it with an emissions trading system (ETS) whose
allowance price hovered around EUR~75 per tonne in early 2026\cite{ecprice}.
\end{concept}

\begin{concept}{Embedded (embodied) emissions.}
The greenhouse gases released while making a product --- the electricity for
the furnace, the coal behind the steel. A tonne of finished steel fasteners
carries roughly 2~tCO$_2$e of embedded emissions, by industry
estimate\cite{specialinsert}.
\end{concept}

\begin{concept}{CBAM --- the EU's Carbon Border Adjustment Mechanism.}
Since 1~January~2026, importers bringing carbon-intensive goods (cement, iron
and steel, aluminium, fertilisers, electricity, hydrogen) into the EU must buy
``CBAM certificates'' for the embedded emissions in those goods\cite{taxud}.
Steel screws and bolts --- customs code CN~7318 --- are covered\cite{carbonchain}.
\end{concept}

\begin{concept}{Default values.}
If nobody proves a product's real emissions, the EU assigns a deliberately
pessimistic ``default value'' --- set at the highest observed intensity and
marked up 10\% in 2026, rising to 30\% by 2028\cite{omm,bpc}. Real, verified
data is almost always far cheaper. This asymmetry is the engine of the entire
problem.
\end{concept}

\begin{concept}{Scope 1 / 2 / 3.}
A firm's own fuel burning (Scope~1), the emissions behind its purchased
electricity (Scope~2), and everything upstream and downstream in its value
chain (Scope~3). When a giant like TSMC works on its Scope~3, the pressure
lands on its small suppliers\cite{tsmc}.
\end{concept}

With those five ideas, the story is simple: \textbf{Taiwan's export SMEs are
being asked, by market forces they cannot see or negotiate with, to produce
data they have no tools to produce.} The rest of this document quantifies that
sentence and proposes what to build about it.
